\documentclass{situationist-preprint}

\title{Cognitive Phenomenology Defended}
\msid{MND-19-a}
\msfamily{MND-19}
\msdiscipline{Philosophy of Mind}
\mstype{Principal manuscript}
\mskeywords{cognitive phenomenology; occurrent thought; phenomenal contrast; self-knowledge; sensory imagery; introspection}

\begin{document}
\maketitle

\begin{abstract}
Is there something it is like to think, over and above what it is like to have the sensory imagery, inner speech and feelings that accompany thinking? Conservatives deny it; liberals affirm it. The debate has largely been conducted by appeal to introspective reports, and the reports conflict. We argue that occurrent thoughts have a distinctive phenomenal character not reducible to sensory imagery, and we argue for it without resting on introspective reports about the presence of such a character. The argument is from self-knowledge. A subject can know immediately and non-inferentially which thought she is occurrently thinking; immediate knowledge of a current conscious state is knowledge fixed by that state's phenomenal character; and the sensory phenomenal character accompanying a thought underdetermines which thought it is. So the phenomenal character that fixes the knowledge is not sensory. We defend the second premise against the interpretive theory of self-knowledge, and we note that the result concerns what thinking is like, not what makes something a thinker.
\end{abstract}

\section{The problem}

That perceptual experiences, bodily sensations and emotions have phenomenal character is uncontested among those who accept phenomenal character at all. Whether thoughts have it is contested. \emph{Liberals} hold that consciously judging, wondering or understanding has a phenomenal character proper to it, which varies with the content of the thought \citep{Strawson1994,Siewert1998,HorganTienson2002,Pitt2004}. \emph{Conservatives} hold that whatever is phenomenal in thinking is sensory: inner speech, visual imagery, feelings of effort or certainty \citep{Prinz2011,TyeWright2011,CarruthersVeillet2011}.

Liberals have relied chiefly on phenomenal contrast. Strawson's example contrasts a monolingual English speaker and a French speaker listening to the same French broadcast: their auditory experiences are alike, but what it is like for the French speaker includes understanding. Conservatives reply that the difference lies in further sensory states, such as imagery and feelings of familiarity, which the French speaker has and the other lacks. Both sides then report what they find on introspection, and the reports differ. MND-20 argues that classification of an experience under theory-laden categories is a domain where introspection is unreliable, and that this dispute is an example. The dispute must therefore be settled by argument, if at all. We offer one, developing a line due to \citet{Pitt2004}.

\section{Framework}

\begin{definition}[Cognitive phenomenology]\label{d:cp}
An occurrent thought has \emph{cognitive phenomenology} if there is something it is like to think it that is not constituted by the phenomenal character of sensory states (perceptual experience, imagery, inner speech, bodily and emotional feeling) that accompany it, and that varies with the content of the thought.
\end{definition}

\begin{principle}[Immediate knowledge of thought]\label{p:know}
A subject who occurrently and consciously thinks a thought is normally in a position to know, immediately and without inference from other evidence, which thought she is thinking.
\end{principle}

\begin{principle}[Phenomenal basis]\label{p:basis}
Immediate, non-inferential knowledge of which current conscious state one is in is knowledge whose content is fixed by the state's phenomenal character.
\end{principle}

Principle~\ref{p:basis} is the account of introspective knowledge defended in MND-20: a judgment about a current phenomenal state is reliable because it can be formed by a process in which the state itself fixes the content of the judgment. What distinguishes immediate from inferential self-knowledge, on that account, is that in the former the state is not evidence for a conclusion but a constituent of what is judged.

\begin{assumption}[Underdetermination]\label{a:under}
The sensory states accompanying an occurrent thought underdetermine which thought it is: the same imagery and inner speech can accompany different thoughts, and different imagery can accompany the same thought.
\end{assumption}

Assumption~\ref{a:under} is accepted by conservatives. Inner speech of the sentence \enquote{visiting relatives can be boring} accompanies two thoughts, one about the activity, one about the relatives; a thought that the meeting is at noon can be accompanied by an image of a clock, by inner speech, or by nothing sensory one can find. The point is about the sensory states themselves, not about what one reports.

\section{The argument}

\begin{theorem}\label{t:main}
Occurrent thoughts have a distinctive phenomenal character not reducible to sensory imagery.
\end{theorem}

\begin{proof}
Let $S$ think the thought that the relatives who are visiting can be boring, with inner speech of the ambiguous sentence and no further imagery. By Principle~\ref{p:know}, $S$ can know immediately that this is the thought she is thinking, and not the thought that the activity of visiting relatives can be boring. By Principle~\ref{p:basis}, this knowledge has its content fixed by the phenomenal character of her current state. By Assumption~\ref{a:under}, the sensory phenomenal character of her state is the same as it would be if she were thinking the other thought; so it cannot fix which of the two she is thinking. Hence her current state has a phenomenal character that is not sensory and that varies with the content of the thought. By Definition~\ref{d:cp}, the thought has cognitive phenomenology.
\end{proof}

The argument does not ask the reader to find, on introspection, a non-sensory phenomenal character. Its premises concern the existence of immediate self-knowledge, its basis, and a fact about sensory states. A conservative who rejects the conclusion must reject one of them.

\section{Objections and replies}

\begin{objection}[Interpretation]\label{o:isa}
Knowledge of one's own thoughts is not immediate. It is reached by the same interpretive mechanisms we use for others, applied to one's own inner speech, imagery and situation \citep{Carruthers2011}. Principle~\ref{p:know} is false, and Theorem~\ref{t:main} has no basis.
\end{objection}

\begin{reply}
The interpretive theory is plausible for knowledge of the causes of one's states and of one's standing attitudes, and MND-20 grants that introspection is unreliable about causes. It is implausible for the occurrent thought. Interpretation of ambiguous inner speech from context would leave a subject uncertain which reading she meant in cases where the context does not discriminate, and subjects are not uncertain in such cases; and the interpretive theory must explain why the result of interpretation is so reliably correct in exactly the cases where the evidence it works on is ambiguous. If the reply is that an unconscious representation of the meant content guides the interpretation, the subject's knowledge is fixed by that content after all, and the remaining question is only whether it is fixed through its conscious presence, which is Principle~\ref{p:basis}.
\end{reply}

\begin{objection}[Non-phenomenal access]\label{o:access}
Immediate knowledge of thought may rest on access to the thought's content without any phenomenal character: thoughts are access-conscious, not phenomenally conscious.
\end{objection}

\begin{reply}
Then immediate knowledge of thought would be the one case of introspective knowledge of a current conscious state not grounded in its phenomenal character, and the asymmetry would need an argument. The objector's best argument is that thoughts lack phenomenal character; but that is the conclusion in dispute. Principle~\ref{p:basis} gives a uniform account of immediate self-knowledge; the objection gives a disjunctive one without independent support.
\end{reply}

\begin{objection}[Hidden imagery]\label{o:hidden}
The two readings of the ambiguous sentence are accompanied by different, perhaps fleeting, images: of relatives sitting in one's house, of a train journey.
\end{objection}

\begin{reply}
Such images can be absent, and when present they too underdetermine content: an image of relatives in one's house is compatible with either reading. Assumption~\ref{a:under} concerns all accompanying sensory states, and the objection can only multiply them, not make them determine content.
\end{reply}

\section{What the result does not establish}

The theorem concerns what thinking is like, not what thinking is. MND-08 argues that thinking is a functional property, so that whether a system thinks is fixed by its functional organization. Theorem~\ref{t:main} is consistent with this: in conscious subjects, occurrent thinking is accompanied by a phenomenal character of its own, but that character is not what makes something a thinker. Nor does the result settle whether cognitive phenomenology is a phenomenology of attitude, of content, or both, or whether it can occur without any accompanying sensory state.

\section{Conclusion}

A subject knows immediately which thought she is thinking; such knowledge is fixed by the phenomenal character of the state known; and no sensory character fixes which thought it is. So occurrent thoughts have a phenomenal character of their own. The dispute over cognitive phenomenology is thereby settled by argument rather than by competing introspective reports.

\begin{thebibliography}{99}\small
\bibitem[Carruthers(2011)]{Carruthers2011} Carruthers, P. (2011). \emph{The Opacity of Mind: An Integrative Theory of Self-Knowledge}. Oxford University Press.
\bibitem[Carruthers and Veillet(2011)]{CarruthersVeillet2011} Carruthers, P. and Veillet, B. (2011). The case against cognitive phenomenology. In T. Bayne and M. Montague (eds.), \emph{Cognitive Phenomenology}. Oxford University Press.
\bibitem[Horgan and Tienson(2002)]{HorganTienson2002} Horgan, T. and Tienson, J. (2002). The intentionality of phenomenology and the phenomenology of intentionality. In D.~J. Chalmers (ed.), \emph{Philosophy of Mind: Classical and Contemporary Readings}. Oxford University Press.
\bibitem[Pitt(2004)]{Pitt2004} Pitt, D. (2004). The phenomenology of cognition, or, what is it like to think that P? \emph{Philosophy and Phenomenological Research}, 69(1), 1--36.
\bibitem[Prinz(2011)]{Prinz2011} Prinz, J. (2011). The sensory basis of cognitive phenomenology. In T. Bayne and M. Montague (eds.), \emph{Cognitive Phenomenology}. Oxford University Press.
\bibitem[Siewert(1998)]{Siewert1998} Siewert, C. (1998). \emph{The Significance of Consciousness}. Princeton University Press.
\bibitem[Strawson(1994)]{Strawson1994} Strawson, G. (1994). \emph{Mental Reality}. MIT Press.
\bibitem[Tye and Wright(2011)]{TyeWright2011} Tye, M. and Wright, B. (2011). Is there a phenomenology of thought? In T. Bayne and M. Montague (eds.), \emph{Cognitive Phenomenology}. Oxford University Press.
\end{thebibliography}
\end{document}
